\section{Stage IV: Strength Reduction in the Compiler}
\label{sec:sr}

\subsection{The Missing Optimisation}
In the inner loop \code{b_k} is invariant and \code{i} grows by a constant, so
\code{b_k + i} is itself an induction variable that could be incremented per
iteration instead of rebuilt with a \code{mov} and an \code{add}: \emph{strength
reduction}~\cite{cocke1977strength}. A six-line reproducer shows that flambda2
does not do it:

\begin{lstlisting}[language=OCamlX,caption={Strength-reduction reproducer.},label=lst:sr-repro,numbers=none]
let f (d : float array) (x : float array) b0 cols =
  let a = ref 0.0 in
  for i = 0 to cols - 1 do a := !a +. (d.(b0 + i) *. x.(i)) done; !a
\end{lstlisting}

\begin{lstlisting}[language=asmX,caption={Its inner loop under OxCaml 5.2.0+ox, \texttt{-O3 -unsafe}.},label=lst:sr-asm]
.L109: movq    %rdi, %rcx                # reload the invariant b0
       addq    %rdx, %rcx                # recompute b0 + i, every iteration
       vmovsd  -8(%rax,%rcx,4), %xmm1    # d.(b0+i): two extra instructions
       vmulsd  -4(%rbx,%rdx,4), %xmm1, %xmm1   # x.(i): none, i is the counter
       vaddsd  %xmm1, %xmm0, %xmm0
       cmpq    %rsi, %rdx
       je      .L118
       addq    $2, %rdx                  # i++ (tagged)
       cmpq    (%r14), %r15              # poll
       jbe     .L126
       jmp     .L109
\end{lstlisting}

In the same loop \code{x.(i)} costs no extra instructions because \code{i} is
the counter, while \code{d.(b0 + i)} costs two, with or without bounds checks.

Is this cost imposed by the language or by the compiler? Writing the C++
kernel in OCaml's index form, \code{d[b0 + i]}, separates the two: if the
index form were inherently costly, it would carry the same extra
instructions; if the cost lies in the compiler, GCC would remove them. It
compiles to the same 15-instruction loop as the pointer version and runs
equally fast (a separate 21-round run): GCC strength-reduces the index. The
language forces OCaml into the index form, but only its compiler leaves it
unoptimised, so a compiler pass is the right fix.

\subsection{By Hand First}
Carrying four running indices \code{j0..j3} in the source (build switch
\code{MG_SR=1}) removed the eight rebuild instructions and the spill (on
5.2.0+ox the loop shrank from 23 to 21 instructions), gaining
@P{oxcaml_unsafe_unroll,oxcaml_unsafe_sr}@ on the OxCaml release.

On stock OCaml 5.3.0 the same edit gave
@P{ocaml_unsafe_unroll,ocaml_unsafe_sr}@. But a source-level workaround must be
repeated by hand in every loop and makes the source harder to read.

\subsection{Confirming the Gap in the Compiler Source}
Neither flambda2's catalogue of optimisations nor any of 29,035 commits
mentions loops, induction or strength; there is no loop-invariant code motion,
value numbering or unrolling. My guide
suggested contacting the OCaml maintainers and first opening the change as a
pull request to my own fork.

\subsection{A Strength-Reduction Pass over Cmm}
\label{sec:sr-pass}
I implemented the pass on Cmm rather than in flambda2
(Figure~\ref{fig:pipeline}). Cmm has no assignment, so ``bound outside the
loop'' \emph{is} the invariance analysis.

\begin{figure}[h]
  \centering
  \includegraphics[width=\textwidth]{diag_pipeline}
  \caption{The new pass in the OxCaml backend; the checks on the right make
  interior pointers safe (Section~\ref{sec:ip}).}
  \label{fig:pipeline}
\end{figure}

The pass finds a counter \code{c} that every back edge steps by a constant
$n$, collects variables affine in it ($m\cdot c + k$), and rewrites every sum
\code{inv + idx}, with \code{idx} affine and \code{inv} an integer bound
outside the loop, into a loop parameter starting at \code{inv + (m*c0 + k)}
and stepping by $m\cdot n$:

\begin{lstlisting}[language=CmmX,caption={The transformation in pseudo-Cmm (tagged index $i = 2c + 1$).},label=lst:sr-cmm]
; before
(catch rec (loop c a0 a1 a2 a3)
   ... (load float64 (+a (+a d (<< (+ b1 i) 2)) -8)) ...
   (exit loop (+ c 1) ...))
; after: b1 + i is carried as sr1
(catch rec (loop c a0 a1 a2 a3 sr1)
   ... (load float64 (+a (+a d (<< sr1 2)) -8)) ...
   (exit loop (+ c 1) ... (+ sr1 2)))
with entry (exit loop 0 0. 0. 0. 0. (+ b1 1))
\end{lstlisting}

Both operands must be integers (a carried value is live across the back-edge
poll, so it must never be a pointer), and variables bound more than once are
ignored; wrapping addition is associative, so overflow behaviour is
unchanged. Three cost rules decide whether a rewrite pays:

\begin{itemize}
  \item \textbf{Register neutrality.} The invariant must have \emph{no other
        use} in the loop, so it dies at entry.
  \item \textbf{Only sums that must be materialised.} OCaml's tagged indices
        reach Cmm as \code{(base + i) - 1}, which x86 forms with one
        \code{lea}; carrying such a sum only adds an \code{add} and a live
        register. Skipping these sums makes nx-oxcaml's element-wise kernels
        1.10--1.14$\times$ faster under default flags, the passes-off speed
        (a separate 15-round run).
  \item \textbf{At most eight carried values.} Without this cap a loop with 24
        such sums ran 1.65$\times$ slower than with the pass off; with it, the
        loop runs at the passes-off speed.
\end{itemize}

The pass sits behind \code{-[no-]cmm-strength-reduction}; with the flag off
the code is byte-identical.

\subsection{Results}
\begin{itemize}
  \item \textbf{Correctness.} 17 hand-written edge cases (counting down,
        overflow, exceptions, forced major GCs) and 30 random programs match
        the pass off; Section~\ref{sec:counted} gives the test suite.
  \item \textbf{Speed on unmodified source.} @P{pr_c_off,pr_c_sr}@ on the
        final compiler with counted loops (Table~\ref{tab:passes}), and
        @P{ox54_base,ox54_sr}@ on the compilers built without polls, close to
        the hand edit's @D{oxcaml_unsafe_unroll,oxcaml_unsafe_sr}@.
  \item \textbf{Code.} The pass leaves one re-tag the hand edit does not,
        \ainl|leaq 1(%r10,%r10)|, from the translation to Cmm; interior
        pointers (Section~\ref{sec:ip}) remove it.
\end{itemize}

The pass does not cover non-constant strides (Section~\ref{sec:ivs}) or
invariant expressions that flambda2 has inlined into an address, which
explains the crypto-study result (Section~\ref{sec:crypto}).
